% Generated draft: requires mathematical and source review.
\documentclass[11pt]{article}
\usepackage[margin=25mm]{geometry}
\usepackage{fontspec}
\setmainfont{Latin Modern Roman}
\usepackage{amsmath,amssymb,mathtools,unicode-math}
\setmathfont{Latin Modern Math}
\usepackage{enumitem,needspace,xurl,hyperref}
\setlength{\parindent}{0pt}
\setlength{\parskip}{5pt}
\setlength{\emergencystretch}{3em}
\newcommand{\N}{\mathbb{N}}
\newcommand{\Z}{\mathbb{Z}}
\newcommand{\Q}{\mathbb{Q}}
\newcommand{\R}{\mathbb{R}}
\newcommand{\C}{\mathbb{C}}
\newcommand{\F}{\mathbb{F}}
\newcommand{\A}{\mathbb{A}}
\newcommand{\PP}{\mathbb{P}}
\newcommand{\E}{\mathbb{E}}
\newcommand{\eps}{\varepsilon}
\newcommand{\dd}{\,\mathrm{d}}
\begin{document}
\section*{UnsolvedMath ID: 30002492}
\textbf{Source problem code:} catalogue entry
\subsection{Definitions and mathematical statement}
Under the standard interpretation of the weight, let \(k\geq 2\) be an even integer, and let
\[
S_k(1)=S_k(\mathrm{SL}_2(\Z))
\]
be the complex vector space of holomorphic cusp forms of weight \(k\) for the full modular group. Write \(m_k=\dim_{\C}S_k(1)\). A normalized Hecke eigenform is a nonzero simultaneous eigenvector for all Hecke operators \(T_n\), with Fourier expansion
\[
f(z)=\sum_{n\geq 1}a_n(f)q^n,\qquad q=e^{2\pi i z},\qquad a_1(f)=1.
\]
Its coefficient field is the number field
\[
\Q_f=\Q\bigl(a_n(f):n\geq 1\bigr)\subseteq\C.
\]
Let \(L_f\subseteq\C\) be the normal closure of \(\Q_f/\Q\), and let \(\mathfrak S_m\) denote the symmetric group on \(m\) letters.

\textbf{Maeda's conjecture, as formulated in the supplied record.} For every such weight \(k\) and every normalized Hecke eigenform \(f\in S_k(1)\), does one have
\[
[\Q_f:\Q]=m_k
\qquad\text{and}\qquad
\operatorname{Gal}(L_f/\Q)\cong\mathfrak S_{m_k}?
\]
When \(S_k(1)=\{0\}\), there is no normalized eigenform, so the assertion at that weight is vacuous. The supplied background also describes the prediction that all normalized Hecke eigenforms of a fixed weight form one Galois orbit, where Galois automorphisms act on Fourier coefficients.

\textit{Formulation qualifications.} The source does not explicitly specify the range of \(k\); positive even integral weights are the standard convention used here and require confirmation against the cited report. The original wording quantifies over every \(f\in S_k(1)\), whereas the supplied cleaned statement restricts to normalized Hecke eigenforms. The formulation above follows that explicit restriction; it does not assert the claim for arbitrary cusp forms.
\subsection{Short English statement}
For each level-one weight, Maeda's conjecture predicts that every normalized cuspidal Hecke eigenform has coefficient-field degree equal to the dimension of the cusp-form space, and that its normal closure has the full symmetric Galois group. The supplied background also predicts a single Galois orbit. The weight convention and the restriction from arbitrary cusp forms to normalized eigenforms are explicit interpretation qualifications.
\subsection{Sources}
\begin{itemize}
\item Local UnsolvedMath registry, supplied complete record: permanent ID \(30002492\), title \emph{Maeda Conjecture for Level-One Cusp Forms}, exact source problem code \texttt{catalogue entry}, in the Oberwolfach Reports Open Problems collection.
\item Supplied report citation: \emph{Langlands Correspondence and Constructive Galois Theory} (2014), Oberwolfach Reports, pp.\ 287--334, DOI 10.4171/owr/2014/05. \url{https://doi.org/10.4171/owr/2014/05}. The record attributes the conjecture to Maeda but supplies no individual proposer for this report entry; 2014 is the supplied report/proposed-year metadata.
\item Supplied literature source: \emph{Galois orbits of newforms}. \url{https://doi.org/10.4171/JEMS/1073}.
\item Supplied literature source: \emph{computational evidence for Maeda's conjecture}. \url{https://arxiv.org/abs/1207.3480}.
\item Supplied literature assessment, checked at \texttt{2026-08-21}: the record labels the problem open and states that the single-orbit, full-degree, and symmetric-normal-closure predictions remain unproved in general, despite extensive computation and partial results for related levels. This is supplied dated curation derived from the OWR curation, not an independently verified assessment; its time-sensitive status and cited claims require independent checking.
\item The supplied statement-status label is \texttt{corrected\_verified}. Its accompanying verification note says that the mathematical content was unchanged and that notation, typography, and wording were normalized into a self-contained publication form. These are supplied editorial claims, not mathematical verification of this draft.
\end{itemize}
\end{document}
